% ============================================================
% chapter4.tex — Chương 4: Kết quả dự kiến và Kế hoạch thực hiện
% ============================================================

\chapter{KẾT QUẢ DỰ KIẾN VÀ KẾ HOẠCH THỰC HIỆN}
\label{chap:kehoach}

% ---------------------------------------------------------------
\section{Sản phẩm dự kiến}
\label{sec:sanpham}
% ---------------------------------------------------------------

Bảng~\ref{tab:deliverables} liệt kê các sản phẩm dự kiến cùng tiêu chí nghiệm thu của đề tài.

\begin{table}[H]
\centering
\caption{Danh sách sản phẩm dự kiến và tiêu chí nghiệm thu}
\label{tab:deliverables}
\begin{tabularx}{\textwidth}{l c X}
\toprule
\textbf{Sản phẩm} & \textbf{Bắt buộc} & \textbf{Tiêu chí nghiệm thu / Đạt} \\
\midrule
Báo cáo nghiên cứu & Có & Đầy đủ cấu trúc đề cương nghiên cứu khoa học, trích dẫn chuẩn hóa. \\
Module BIM/IFC-to-Graph & Có & Trích xuất thành công node/edge từ file \texttt{.ifc} và xuất ra tệp tin graph JSON. \\
Building graph tòa nhà & Có & Biểu diễn đúng topology phòng/hành lang/cầu thang dưới dạng NetworkX graph. \\
Bộ kịch bản cháy CFAST & Có & Hàng trăm kịch bản cháy được sinh tự động bằng script Python. \\
FDS reference cases & Có & Tối thiểu 3 case FDS P0 được thiết lập, căn chỉnh sensor schema và chạy thành công để đối sánh với CFAST. \\
Virtual IoT dataset & Có & Sinh được chuỗi thời gian dữ liệu cảm biến ảo 3 lớp (clean, noisy, faulty) từ mô phỏng. \\
Baseline models & Có & Có rule-based, ít nhất 1 baseline không graph và ít nhất 1 baseline graph, dùng cùng split và có metric P0. \\
Temporal GNN/BAT-GNN & Có & Dự báo được \texttt{smoke\_label} tại $t+\Delta$ và \texttt{time\_to\_arrival}, có đối sánh với baseline mạnh nhất. \\
Dashboard prototype & Có & Replay dữ liệu pipeline và hiển thị hai đầu ra P0, không dùng số liệu hard-code. \\
Evaluation report & Có & Có scenario/fire-node/fire-floor split, FDS reference set, ablation, robustness, inference time và phân tích giới hạn. \\
\bottomrule
\end{tabularx}
\end{table}

% ---------------------------------------------------------------
\section{Kế hoạch thực hiện}
\label{sec:kehoachthuchien}
% ---------------------------------------------------------------

Đề tài được thực hiện trong 20 tuần, bắt đầu từ ngày 19/06/2026. Kế hoạch chi tiết được trình bày trong Bảng~\ref{tab:roadmap} và biểu đồ Gantt (Hình~\ref{fig:gantt}).

\begin{table}[H]
\centering
\caption{Lộ trình thực hiện 20 tuần}
\label{tab:roadmap}
\begin{tabularx}{\textwidth}{c X X}
\toprule
\textbf{Tuần} & \textbf{Công việc chính} & \textbf{Sản phẩm} \\
\midrule
1  & Chốt phạm vi, tạo repo, ghi chú giả thiết         & Tài liệu scope \\
2  & Export IFC, kiểm tra mô hình BIM                    & File IFC + ghi chú \\
3  & Parse room/door/window/floor/stair                  & Bảng dữ liệu thô \\
4  & Xây dựng và trực quan hóa graph                    & Graph files \\
5  & Tạo case CFAST debug (layout thấp tầng)             & Case CFAST đầu tiên \\
6  & Chạy FDS pilot và chốt quy tắc căn chỉnh           & Báo cáo pilot CFAST--FDS \\
7  & Schema kịch bản + batch runner                      & Scenario generator \\
8  & Parse CFAST, tạo nhãn P0 và batch dataset           & Parsed outputs + labels \\
9  & Tạo virtual sensor clean                            & Clean stream \\
10 & Tạo noisy/faulty IoT presets                        & IoT streams \\
11 & Rule-based + baseline không graph                   & Baseline metrics \\
12 & Baseline graph                                      & Bảng so sánh \\
13 & Temporal GNN / BAT-GNN v1                           & Trained model \\
14 & Chuẩn bị các FDS reference case P0 còn lại          & FDS input \\
15 & Chạy tối thiểu 3 FDS case P0                        & FDS raw output \\
16 & Tạo sensor/label FDS và đối sánh                    & Báo cáo CFAST--FDS \\
17 & Ablation + robustness study                         & Figures đánh giá \\
18 & Dashboard replay + inference API                    & Demo \\
19 & Viết báo cáo draft                                  & Bản thảo báo cáo \\
20 & Hoàn thiện slides, demo, nộp                        & Bài nộp cuối \\
\bottomrule
\end{tabularx}
\end{table}

\begin{figure}[H]
\centering
\resizebox{0.95\textwidth}{!}{%
\begin{ganttchart}[
  x unit=0.45cm,
  y unit title=0.45cm,
  y unit chart=0.35cm,
  hgrid,
  vgrid,
  title height=1,
  bar height=0.5,
  bar/.append style={fill=NavyBlue!60},
  bar label font=\tiny,
  title label font=\tiny\bfseries,
  milestone label font=\tiny,
  group label font=\tiny\bfseries,
  group/.append style={fill=NavyBlue!30, draw=NavyBlue},
]{1}{20}
  \gantttitle{Tuần}{20} \\
  \gantttitlelist{1,...,20}{1} \\

  % BIM/Graph
  \ganttgroup{BIM \& Graph}{1}{4} \\
  \ganttbar{Scope \& repo}{1}{1} \\
  \ganttbar{Export IFC}{2}{2} \\
  \ganttbar{Parse IFC}{3}{3} \\
  \ganttbar{Build graph}{4}{4} \\

  % Simulation
  \ganttgroup{Mô phỏng}{5}{8} \\
  \ganttbar{CFAST debug}{5}{5} \\
  \ganttbar{FDS pilot}{6}{6} \\
  \ganttbar{Scenario batch}{7}{7} \\
  \ganttbar{Parse + labels}{8}{8} \\

  % IoT
  \ganttgroup{IoT}{9}{10} \\
  \ganttbar{Clean sensors}{9}{9} \\
  \ganttbar{Noisy/faulty}{10}{10} \\

  % AI
  \ganttgroup{AI Models}{11}{13} \\
  \ganttbar{Baselines}{11}{12} \\
  \ganttbar{Temporal GNN}{13}{13} \\

  % FDS
  \ganttgroup{FDS Reference}{14}{16} \\
  \ganttbar{Prep + Run FDS}{14}{15} \\
  \ganttbar{Sensors + Compare}{16}{16} \\

  % Evaluation
  \ganttgroup{Đánh giá \& Demo}{17}{18} \\
  \ganttbar{Ablation}{17}{17} \\
  \ganttbar{Dashboard}{18}{18} \\

  % Report
  \ganttgroup{Báo cáo}{19}{20} \\
  \ganttbar{Draft}{19}{19} \\
  \ganttbar{Submit}{20}{20} \\

\end{ganttchart}
}%
\caption{Biểu đồ Gantt kế hoạch thực hiện 20 tuần}
\label{fig:gantt}
\end{figure}

% ---------------------------------------------------------------
\section{Phân công nhiệm vụ}
\label{sec:phancong}
% ---------------------------------------------------------------

\begin{table}[H]
\centering
\caption{Phân công nhiệm vụ}
\label{tab:phancong}
\begin{tabularx}{\textwidth}{l X l}
\toprule
\textbf{Thành viên} & \textbf{Nhiệm vụ chính} & \textbf{Tuần} \\
\midrule
\svA & Quản lý chung, BIM/IFC-to-Graph, báo cáo        & 1--4, 19--20 \\
\svD & Mô phỏng CFAST, post-processing, labels          & 5--8 \\
\svB & Virtual IoT, dataset builder                      & 9--10 \\
\svC & AI models (baselines + Temporal GNN + BAT-GNN)    & 11--13, 17 \\
\svE & FDS pilot/đối sánh, dashboard prototype           & 6, 14--16, 18 \\
\bottomrule
\end{tabularx}
\end{table}

\begin{notebox}
Tất cả thành viên đều tham gia viết báo cáo ở tuần 19--20. Phân công có thể được điều chỉnh trong quá trình thực hiện.
\end{notebox}

% ---------------------------------------------------------------
\section{Điều kiện thực hiện}
\label{sec:dieuken}
% ---------------------------------------------------------------

\subsection{Phần mềm}

\begin{itemize}
  \item \textbf{BIM:} Autodesk Revit (export IFC).
  \item \textbf{Mô phỏng:} CFAST (NIST), FDS (NIST), Smokeview.
  \item \textbf{Lập trình:} Python 3.10+, PyTorch, PyTorch Geometric, IfcOpenShell, NetworkX, Pandas.
  \item \textbf{Dashboard:} FastAPI, React, TypeScript, Cytoscape.js.
  \item \textbf{Quản lý:} Git, Docker Compose.
  \item \textbf{Báo cáo:} \LaTeX\ (XeLaTeX).
\end{itemize}

\subsection{Phần cứng}

\begin{itemize}
  \item Máy tính cá nhân với GPU hỗ trợ CUDA (cho huấn luyện AI).
  \item Có thể sử dụng Google Colab hoặc cloud GPU nếu cần.
  \item FDS yêu cầu CPU đa nhân; thời gian chạy phụ thuộc quy mô mesh.
\end{itemize}

\subsection{Tài liệu và dữ liệu}

\begin{itemize}
  \item File Revit tòa nhà thực tế (do nhóm chuẩn bị hoặc thu thập).
  \item Tài liệu kỹ thuật CFAST, FDS, IFC từ NIST và buildingSMART.
  \item Quy chuẩn QCVN 06:2022/BXD và các tiêu chuẩn liên quan.
  \item Các bài báo khoa học liên quan về GNN, fire prediction, BIM-to-simulation.
\end{itemize}

% ---------------------------------------------------------------
\section{Rủi ro và phương án giảm thiểu rủi ro}
\label{sec:ruiro}
% ---------------------------------------------------------------

Trong quá trình thực hiện đề tài, nhóm xác định được 10 rủi ro kỹ thuật chính và đề xuất các phương án giảm thiểu tương ứng được trình bày trong Bảng~\ref{tab:ruiro}.

\begin{table}[htbp]
\centering
\caption{Bảng phân tích rủi ro và phương án giảm thiểu}
\label{tab:ruiro}
\footnotesize
\begin{tabularx}{\textwidth}{
  c
  >{\raggedright\arraybackslash}p{0.18\textwidth}
  >{\raggedright\arraybackslash}p{0.28\textwidth}
  >{\raggedright\arraybackslash}X}
\toprule
\textbf{STT} & \textbf{Yếu tố rủi ro} & \textbf{Hệ quả tiềm ẩn} & \textbf{Phương án fallback / Giảm thiểu} \\
\midrule
1 & Tài liệu không xác minh được & Lỗi học thuật nghiêm trọng & Kiểm tra metadata theo DOI hoặc trang chính thức của nhà xuất bản/cơ quan phát hành trước khi trích dẫn. \\
2 & Không có dữ liệu cháy thực tế & AI học mẫu giả lập, sai khác thực tế & Ghi rõ giới hạn; không tuyên bố thay thế CFD thật; chỉ dùng làm surrogate model nghiên cứu. \\
3 & Ngưỡng hazard thiếu chuẩn hóa & Nhãn vùng nguy hiểm không đáng tin & Giữ hazard ở mức P1; chỉ triển khai khi xác minh được ngưỡng và luôn ghi rõ đây không phải kết luận pháp lý. \\
4 & File IFC bị lỗi khi export & Pipeline dữ liệu bị gián đoạn từ đầu & Sử dụng layout thấp tầng đơn giản để debug; vẽ lại thủ công kết nối thiếu trong NetworkX. \\
5 & Mô phỏng FDS quá nặng & Trễ tiến độ đối sánh của đề tài & Chạy pilot ở tuần 6, giảm quy mô mesh và chỉ cam kết tối thiểu 3 case P0. \\
\bottomrule
\end{tabularx}
\end{table}

\begin{table}[htbp]
\centering
\caption*{Bảng~\ref{tab:ruiro} (tiếp theo)}
\footnotesize
\begin{tabularx}{\textwidth}{
  c
  >{\raggedright\arraybackslash}p{0.18\textwidth}
  >{\raggedright\arraybackslash}p{0.28\textwidth}
  >{\raggedright\arraybackslash}X}
\toprule
\textbf{STT} & \textbf{Yếu tố rủi ro} & \textbf{Hệ quả tiềm ẩn} & \textbf{Phương án fallback / Giảm thiểu} \\
\midrule
6 & Rò rỉ dữ liệu (data leakage) & Đánh giá sai hiệu năng AI (metric ảo) & Áp dụng quy tắc split độc lập theo kịch bản cháy; không chia chéo dữ liệu chuỗi thời gian. \\
7 & Mô hình AI bị overfit & Khả năng tổng quát hóa kém trong case study & Áp dụng dropout, weight decay, early stopping, hold-out fire node và hold-out fire floor; không tuyên bố tổng quát hóa sang tòa nhà mới. \\
8 & Luồng IoT giả lập quá sạch & AI không hoạt động được trong thực tế & Xây dựng lớp \texttt{iot\_noisy\_stream} bổ sung nhiễu Gauss, trễ cảm biến và lỗi mất dữ liệu ngẫu nhiên. \\
9 & Dashboard bị hard-code & Demo rỗng, không chạy thực tế được & Kết nối API Backend FastAPI để replay dữ liệu động từ tệp mô phỏng và chạy inference trực tiếp. \\
10 & Bùng nổ tổ hợp kịch bản cháy & Vượt quá khả năng tính toán & Áp dụng phương pháp scenario sampling có kiểm soát thay vì thử nghiệm toàn bộ Cartesian product. \\
\bottomrule
\end{tabularx}
\end{table}
